\cotefile{87}
\scanpage{15}
% L1–3. Vertical rule at the left of the assertion.
\uline{Cor. 6} Soit $F\in\mathrm{Fac}(C)$ ($C$ polyèdre convexe)
Alors pour tt $-1\leq i\leq d=\dim F$, existe $F'\leq F$
telle que $\dim F'=i$.

% L4–9. Horizontal strokes at the two ends, oblique strokes across the block.
\deleted{Clair si $d=-1$ ou $0$, sinon. On \unsure{ad.|dit} qu'il
\inserted{\unsure{procède} par} récurrence sur $d$, clair si $d=-1,0$. Si $d\unsure{\geq|>}0$,
$\exists\ F'$ tel que $F'\ne\varnothing,F$\quad $\varnothing<F'<F$ (cor. 4)
Donc $\dim F'=d'<d$. Si $i\leq d'$ on gagne, par
hyp. de récurrence appliquée à $F'$. \unsure{Si}
Sinon on gagne par l'hyp. de récurrence.}

% L10–15
Formel à partir du cor. 5, qui implique que
si $x\unsure{\leq|<} y$, $x,y\in\mathrm{Fac}(C)$, \deleted{\unsure{alors}} $\exists$ chaîne
$x_0=x<x_1<\ldots x_n=y$, avec $x_{i+1}$ succ de $x_i$
$(0\leq i\leq n-1)$, et qu'alors \deleted{$d(y)-d(x)=n$.}
\deleted{$d\unsure{(x)|(}\quad d\unsure{(}$} $d(x_i)-d(x_0)=i$ i.e. $d(x_i)=d(x_0)+i$.
[On \unsure{prendra|prend} $y=F$, $x=\varnothing$ \ldots{} ]

% L16–20. Vertical rule at the left of the assertion.
\uline{Cor. 7} Soit $\unsure{P|\mathcal P|\Phi}$ un polyèdre combinatoire
strict, $d$ une fonction dimension sur $P$
(cf cor. 5) \unsure{Soient|Soit} $x,y\in P$, $x<y$,
$d(y)=d(x)+2$. Alors l'ens des $z\in P$,
tels que $x<z<y$ est de cardinal $2$.

% L21–24
$P_{x,y}$ est un polyèdre combinatoire \inserted{strict} de
dim $(-1)+2=1$, donc provient d'un polyèdre
convexe de dim $1$, i.e. d'un segment. Celui-ci a exactement $2$ extrémités !
% archive sheet number: 14
